\chapter{TINJAUAN PUSTAKA}
\label{bab:tinjauan_pustaka}

\section{Landasan Teori}
\label{sec:landasan_teori}

Dummy konten landasan teori. Bagian ini menjelaskan konsep-konsep fundamental yang relevan dengan penelitian.

\subsection{Sub-bagian Landasan Teori}
\label{subsec:subbagian_teori}

Deskripsi teori pendukung disertai kutipan yang relevan \cite{dummy2024reference}.

\begin{equation}
    E = mc^2
    \label{eq:einstein}
\end{equation}

Persamaan \ref{eq:einstein} menunjukkan hubungan antara energi dan massa.

\section{Penelitian Terkait}
\label{sec:penelitian_terkait}

Dummy konten penelitian terkait. Bagian ini menguraikan penelitian-penelitian sebelumnya yang relevan.

\subsection{Penelitian Terdahulu}
\label{subsec:penelitian_terdahulu}

Tabel \ref{tab:contoh_tabel} menunjukkan contoh tabel penelitian terdahulu.

\begin{table}[h]
    \centering
    \caption{Contoh Tabel Penelitian Terdahulu}
    \label{tab:contoh_tabel}
    \begin{tabularx}{\textwidth}{X c c c}
        \toprule
        \textbf{Peneliti} & \textbf{Tahun} & \textbf{Metode} & \textbf{Akurasi} \\
        \midrule
        Peneliti A & 2020 & Metode X & 85.0\% \\
        Peneliti B & 2021 & Metode Y & 87.5\% \\
        Peneliti C & 2022 & Metode Z & 90.2\% \\
        \bottomrule
    \end{tabularx}
\end{table}

\section{Kesenjangan Penelitian}
\label{sec:kesenjangan}

Dummy konten kesenjangan penelitian. Bagian ini mengidentifikasi celah yang dapat diisi oleh penelitian ini.

\section{Kerangka Berpikir}
\label{sec:kerangka_berpikir}

Dummy konten kerangka berpikir. Bagian ini menjelaskan alur logis penelitian dari masalah hingga solusi.

Gambar \ref{fig:contoh_gambar} menunjukkan contoh diagram kerangka berpikir.

\begin{figure}[h]
    \centering
    % Gambar placeholder
    \fbox{\parbox[c][4cm][c]{10cm}{\centering Gambar Placeholder}}
    \caption{Contoh Gambar Kerangka Berpikir}
    \label{fig:contoh_gambar}
\end{figure}

\section{Hipotesis}
\label{sec:hipotesis}

Dummy hipotesis penelitian. Berdasarkan tinjauan pustaka, hipotesis penelitian ini adalah:

\begin{enumerate}
    \item Hipotesis pertama.
    \item Hipotesis kedua.
\end{enumerate}